% ============================================================
% HOW TO USE THIS BOOK
% front_matter/how_to_use.tex
% ============================================================

\chapter*{How to Use This Book}
\addcontentsline{toc}{chapter}{How to Use This Book}
\markboth{How to Use This Book}{How to Use This Book}

\textit{The Physics Codex} is designed to be read sequentially.
Each chapter builds on the understanding established in the
chapters before it. Skipping ahead is possible, but you will
learn most effectively by moving through the book in order,
from Chapter 1 to Chapter 29.

\vspace{0.5cm}
\textbf{The structure of each chapter}
\vspace{0.3cm}

Every chapter in this book follows the same structure:

\begin{enumerate}
  \item \textbf{Learning Objectives} --- A brief list of
  what you will be able to do by the end of the chapter.
  Read these before you begin and return to them when you
  finish to check your progress.

  \item \textbf{Opening Quote} --- A statement from a
  scientist, philosopher, or thinker that captures
  something essential about the chapter's subject.
  Think about it.

  \item \textbf{Conceptual Introduction} --- An intuitive,
  everyday explanation of the chapter's main ideas before
  any mathematics is introduced. Read this carefully.
  If you understand this section, the mathematics that
  follows will make complete sense.

  \item \textbf{Main Content} --- The full explanation of
  each topic, including definitions, laws, derivations,
  diagrams, and historical context.

  \item \textbf{Worked Examples} --- Original problems
  solved in full, step by step. \important{Read the entire
  worked example before attempting the exercises. Do not
  look at the solution until you have made a genuine
  attempt.}

  \item \textbf{Chapter Summary} --- A concise review of
  every key concept, formula, and definition in the chapter.
  Use this for revision.

  \item \textbf{Exercises} --- Original practice questions
  divided into two sections:
  \begin{itemize}
    \item \textbf{Section A:} Ten multiple-choice questions
    testing conceptual understanding.
    \item \textbf{Section B:} Five theory or calculation
    questions requiring full written responses.
  \end{itemize}
  Attempt all exercises after reading the chapter and
  studying the worked examples.
\end{enumerate}

\vspace{0.5cm}
\textbf{The box types and what they mean}
\vspace{0.3cm}

You will encounter several types of coloured boxes throughout
the book. Each type serves a specific purpose:

\begin{itemize}
  \item \textcolor{codexblue}{\textbf{Definition boxes}} ---
  Precise, formal definitions of key terms. Memorise these.

  \item \textcolor{codexred}{\textbf{Theorem/Law boxes}} ---
  The fundamental laws and theorems of Physics. Understand
  where they come from, not just what they say.

  \item \textcolor{codexgreen}{\textbf{Worked Example boxes}}
  --- Original solved problems. Study the method, not just
  the answer.

  \item \textcolor{codexgold}{\textbf{Key Concept boxes}} ---
  The most important idea on a given page. If you remember
  nothing else, remember this.

  \item \textcolor{codexgreen!70!black}{\textbf{Real World
  boxes}} --- Connections between the Physics you are learning
  and phenomena you can observe in everyday life. These are
  not optional reading --- they are how you build lasting
  understanding.

  \item \textcolor{codexgold!70!black}{\textbf{Historical
  Note boxes}} --- The human story behind the Physics.
  Context makes concepts memorable.

  \item \textcolor{codexred!80!black}{\textbf{Misconception
  boxes}} --- Common errors students make. Read these
  carefully so you do not make the same mistakes.

  \item \textcolor{codexgold}{\textbf{Exam Tip boxes}} ---
  Specific guidance for UTME and university examinations.
  These contain the kind of strategic advice that separates
  students who understand from students who also perform.
\end{itemize}

\vspace{0.5cm}
\textbf{How to use the answers section}
\vspace{0.3cm}

Answers to all exercises are located at the back of the book,
after the final chapter and appendices. Each answer is labelled
by chapter and question number so you can find it immediately.
For calculation questions, full working is shown --- not just
the final answer. Use the answers to check your work, not to
skip it.

\vspace{0.5cm}
\textbf{A word on mathematics}
\vspace{0.3cm}

Physics requires mathematics. You will encounter algebra,
trigonometry, and basic calculus in this book. If you find
a mathematical step unclear, do not abandon the chapter ---
the physical concept behind the mathematics is always
explained in words as well. The mathematics is a tool.
The understanding is the goal.

\vspace{0.5cm}
\textbf{The appendices}
\vspace{0.3cm}

At the back of the book you will find three appendices:
a complete table of physical constants, a full formula
reference sheet, and a guide to SI units and prefixes.
These are reference tools --- use them when you need them
during study and revision.
